% Sources: README.md; app/backend/newsqa_app/services/chat_service.py; evaluation/testset.py.
\section{Giới thiệu đề tài}
\subsection{Bối cảnh và bài toán}
Tin tức thường phân tán giữa nhiều bài viết có tên người, địa điểm và sự kiện tương tự. Công cụ tìm kiếm trả về danh sách tài liệu nhưng người đọc vẫn phải đối chiếu các đoạn liên quan; một câu trả lời tự nhiên lại khó kiểm chứng nếu không chỉ ra nguồn. NewsLens kết hợp hai nhu cầu này: truy xuất bằng chứng và sinh câu trả lời có dẫn nguồn.

NewsQA cung cấp câu hỏi và đáp án là đoạn trích trong bài CNN~\cite{newsqa}; nhóm chuyển dữ liệu này thành bài toán truy xuất trên toàn kho.
\begin{itemize}
\item \textbf{Đầu vào:} câu hỏi của người dùng và kho bài báo đã chuẩn bị; bài báo và câu hỏi đánh giá đều bằng tiếng Anh.
\item \textbf{Đầu ra RAG:} câu trả lời, chỉ dấu nguồn và đoạn văn hỗ trợ để kiểm tra.
\item \textbf{Người dùng:} người đọc cần tra cứu tin tức và thành viên nhóm cần khảo sát tác động của từng thành phần hệ thống.
\end{itemize}
\subsection{Mục tiêu, câu hỏi nghiên cứu và phạm vi}
Đồ án xây dựng quy trình có thể tái lập và phân biệt lỗi truy xuất với lỗi sử dụng bằng chứng. Bốn câu hỏi nghiên cứu chính là:
\begin{enumerate}
\item Bộ truy xuất và xếp hạng lại nào tìm bằng chứng tốt nhất trong các phương án khảo sát?
\item Lời nhắc zero-shot và số đoạn ngữ cảnh ảnh hưởng thế nào đến chất lượng câu trả lời?
\item Một ví dụ one-shot có cải thiện chất lượng so với zero-shot không?
\item Thay mô hình sinh ảnh hưởng thế nào đến chất lượng và tài nguyên?
\end{enumerate}

Hai khảo sát bổ sung xem xét chiến lược chia đoạn và chính sách từ chối khi thiếu bằng chứng.

Phạm vi gồm chuẩn bị dữ liệu, các vòng thực nghiệm và ứng dụng web sử dụng các thành phần xử lý dùng chung. Kho CNN cố định phục vụ kiểm soát thực nghiệm; chức năng thu thập tin là một nhánh nhập liệu riêng. Đồ án không huấn luyện lại mô hình nền. Chất lượng câu trả lời được đánh giá qua đáp án, bằng chứng và dẫn nguồn thay vì chỉ dựa vào sự xuất hiện của ký hiệu trích dẫn.
\subsection{Đóng góp kỹ thuật}
Các sản phẩm kỹ thuật gồm dữ liệu có nhãn bằng chứng và lịch sử rà soát, các mô-đun truy xuất, xếp hạng lại và sinh câu trả lời, cùng API và giao diện hiển thị nguồn. Kết quả truy xuất được lưu để so sánh lời nhắc trong cùng điều kiện. Cấu hình, dấu định danh và trạng thái từng câu hỏi giúp kiểm tra và tái lập các lần chạy đánh giá.
